\documentclass[10pt,a4paper]{amsart}
\usepackage[margin=19mm]{geometry}
\usepackage{amsmath,amssymb,mathtools}
\usepackage[scr=rsfs,cal=euler]{mathalfa}
\usepackage{xfrac}
\usepackage[unicode,hidelinks,bookmarks=false]{hyperref}
\newcommand{\ie}{i.e.\ }
\newcommand{\cf}{cf.\ }
\newcommand{\resp}{respectively\ }
\newcommand{\Subsection}{\subsection}
\providecommand{\nobreakdash}{\nobreak\hbox{-}}
\setlength{\emergencystretch}{2em}
\multlinegap0pt
\usepackage{graphics}
\usepackage{xy}
\newtheorem{Theorem}{Theorem}[section]
\newtheorem{Lemma}[Theorem]{Lemma}
\newtheorem{Proposition}[Theorem]{Proposition}
\newtheorem{Corollary}[Theorem]{Corollary}
\newtheorem{Definition}[Theorem]{Definition}
\newtheorem{Example}[Theorem]{Example}
\newtheorem{Problem}[Theorem]{Problem}
\newtheorem{Question}[Theorem]{Question}
\newtheorem{Remark}[Theorem]{Remark}
\def\labelenumi{\rm ({\makebox[0.8em][c]{\theenumi}})}
\def\theenumi{\arabic{enumi}}
\def\labelenumii{\rm {\makebox[0.8em][c]{(\theenumii)}}}
\def\theenumii{\roman{enumii}}
\def\@thmcountersep{-}
\def\prend{\hspace*{\fill}\hbox{\rule[-2pt]{3pt}{6pt}}}
\newcommand{\bsquare}{\hbox{\rule{6pt}{6pt}}}
\numberwithin{equation}{section}
\newcommand{\abs}[1]{\lvert#1\rvert}
\begin{document}
\title[Intrinsic linking of simplicial n-complexes in mathbbR^2n: A]{Intrinsic linking of simplicial $n$-complexes in $\mathbb{R}^{2n}$: An additional minimal $n$-complex}
\author{Secondary 57M15 ạte; Ryo Nikkuni}
\date{}
\maketitle
\noindent\emph{Source and licence.} Intrinsic linking of simplicial $n$-complexes in $\mathbb{R}^{2n}$: An additional minimal $n$-complex. Version arXiv:2605.28496v2 (CC BY 4.0). This material is distributed under the Creative Commons Attribution 4.0 International licence, http://creativecommons.org/licenses/by/4.0/, as recorded in the arXiv OAI licence metadata for this exact version and shown on the abstract page https://arxiv.org/abs/2605.28496v2. Full rights notice: licenses/arxiv-2605.28496-CC-BY-4.0.txt.\par\smallskip
\noindent\textit{Editorial scope.} Counted unit: the original \texttt{Theorem} environment \texttt{vko} at source lines 240--245 together with its complete proof at lines 247--255; the delivered document keeps source lines 155--256 so that every referenced label and every citation resolves inside it. Native editable source is the paper's own concatenated TeX; the AMS wrapper is editorial formatting only. No publisher class, upstream script or unrelated artwork is redistributed.
\begin{Theorem}\label{ILR2n} 
Let $n$ be a positive integer. For every embedding $f \colon M^{(n)} \to \mathbb{R}^{2n}$, there exists an element $\lambda \in \Lambda^{n-1,n}(M^{(n)})$ such that ${\rm lk}_2(f(\lambda)) = 1$. 
\end{Theorem}

Note that $K_6$ is isomorphic to $\sigma_5^1$ and that $S(\sigma_5^1)$ contains $M^{(2)}$ as a proper subcomplex. In general, for any positive integer $n$, the suspension $S(\sigma_{2n+1}^{n-1})$ contains exactly $2n+2$ subcomplexes isomorphic to $M^{(n)}$, and each element of $\Lambda^{n-1,n}(S(\sigma_{2n+1}^{n-1}))$ is shared by exactly $n+1$ of these subcomplexes. Thus, for any embedding of $S(\sigma_{2n+1}^{n-1})$ into $\mathbb{R}^{2n}$, Theorem \ref{ILR2n} implies that there exist at least $(2n+2)/(n+1) = 2$ two-component links with ${\rm lk}_2 = 1$. We will prove Theorem \ref{ILR2n} in the next section. While the method in \cite{HRJV26} is based on the original method of Conway--Gordon \cite{CG83}, our proof of Theorem \ref{ILR2n} follows an approach similar to that in \cite{N26b} and is essentially based on the \textit{van Kampen obstruction} \cite{VK33}.




By Theorem \ref{ILR2n}, $S(\sigma_{2n+1}^{n-1})$ is not minimal with respect to this intrinsic linking property. On the other hand, for every positive integer $n$, $M^{(n)}$ is minimal in the following sense.

\begin{Proposition}\label{ILRn2min}
Let $N$ be a proper subcomplex of $M^{(n)}$. Then there exists an embedding $f \colon N \to \mathbb{R}^{2n}$ such that ${\rm lk}_2(f(\lambda)) = 0$ for every element $\lambda \in \Lambda^{n-1,n}(N)$.
\end{Proposition}

Note that $M^{(1)}$ is isomorphic to $N_{2}^{(1)} = K_1 \sqcup K_{2,3}$. On the other hand, the simplicial $n$-complexes $N_{1}^{(n)}$ and $N_3^{(n)}$ each have $2n+3$ vertices, while $N_2^{(n)}$ has $3n+3$ vertices \cite{N26b}. Since $M^{(n)}$ has $2n+4$ vertices, $M^{(n)}$ is not isomorphic to any of the simplicial $n$-complexes $N_{i}^{(n)}$ ($i=1,2,3$) for $n \ge 2$. We will also give a proof of Proposition \ref{ILRn2min} in the next section.


\begin{Remark}\label{k33}
Let $[3]^{*m}$ denote the $m$-fold join of three points. Note that $[3]^{*m}$ can naturally be regarded as a simplicial $(m-1)$-complex. For $n \ge 2$, it is also well known that $[3]^{*n}$ is a minimal simplicial complex that cannot be embedded in $\mathbb{R}^{2n-2}$ \cite{VK33}, \cite{Flores32}. In particular, $[3]^{*2}$ is isomorphic to $K_{3,3}$. In this case, $M_{[3]^{*m}}^{(n)}$ is isomorphic to $N_2^{(n)}$, and the same conclusion as in Theorem~\ref{ILR2n} holds. 
\end{Remark}




\section{The mod $2$ van Kampen obstruction and proof of Theorem \ref{ILR2n}}

Let $K$ be a simplicial $n$-complex. It is well known that every such $K$ admits a generic immersion into $\mathbb{R}^{2n}$. Here, an immersion $\varphi \colon K \to \mathbb{R}^{2n}$ is said to be \textit{generic} if all singularities of $\varphi(K)$ are transversal double points occurring between the interiors of pairs of disjoint $n$-simplices. For a generic immersion $\varphi$ of $K$ into $\mathbb{R}^{2n}$ and any two mutually disjoint $n$-simplices $s, s' \in \varDelta^n(K)$, let $l(\varphi(s),\varphi(s'))$ denote the number of double points between $\varphi(s)$ and $\varphi(s')$. Denote by $K^*$ the \emph{deleted product} of $K$, that is, the subcomplex of $K \times K$ consisting of all cells $s \times s'$ with $s \cap s' = \emptyset$. Let $\widetilde{K}^*$ be the quotient complex of $K^*$ obtained by identifying $s \times s'$ with $s' \times s$. Given a generic immersion $\varphi \colon K \to \mathbb{R}^{2n}$, the modulo-$2$ reduction of $l(\varphi(s),\varphi(s'))$ defines a $\mathbb{Z}_2$-cocycle $c_{\varphi}$ on $\widetilde{K}^*$, and its cohomology class $o_K = [c_{\varphi}] \in H^{2n}(\widetilde{K}^*;\mathbb{Z}_2)$ is called the \emph{mod $2$ van Kampen obstruction} of $K$.



Let $n$ be a positive integer, and let $a$, $b$, and $c$ be three vertices not belonging to $\sigma_{2n}^{n-1}$. In the following, we regard $\sigma_{2n}^{n-1} * \{a,b,c\}$ as a simplicial $n$-complex $K$.

\begin{Lemma}\label{cohomo}
Let $K = \sigma_{2n}^{n-1} * \{a,b,c\}$. Then 
\[
H^{2n}(\widetilde{K}^*;\mathbb{Z}_2) \cong \mathbb{Z}_2.
\]
Moreover, this group is generated by the cohomology class $[c]$ of the cocycle $c$ defined as follows. Fix any two disjoint $(n-1)$-simplices $\tau, \tau' \in \varDelta^{n-1}(\sigma_{2n}^{n-1})$, and let $\sigma = a*\tau$, $\sigma' = b*\tau'$. Define $c \in C^{2n}(\widetilde{K}^*;\mathbb{Z}_2)$ to be the dual cochain of the $2n$-cell $\sigma \times \sigma'$; that is,
\[
c(s \times s') = 
\begin{cases}
1 & \text{if } s \times s' = \sigma \times \sigma' \text{ in } \widetilde{K}^*, \\
0 & \text{otherwise}.
\end{cases}
\]
\end{Lemma}





\begin{proof}
For a pair of disjoint $n$-simplices $s,s'$ in $K$, let $E^{s,s'}$ denote the dual cochain of the $2n$-cell $s \times s'$ in $\widetilde{K}^*$. Similarly, for an $(n-1)$-simplex $t$ and an $n$-simplex $s'$ in $K$ with $t \cap s' = \emptyset$, let $V^{t,s'}$ denote the dual cochain of the $(2n-1)$-cell $t \times s'$ in $\widetilde{K}^*$.

First, fix an $(n-1)$-simplex $t = |a a_0 \cdots a_{n-2}|$ and an $n$-simplex $s' = |b a_{n+1} \cdots a_{2n}|$ in $K$. Note that $t$ and $s'$ are disjoint. Let 
\[
s_1 = |a a_0 \cdots a_{n-2} a_{n-1}|, \quad s_2 = |a a_0 \cdots a_{n-2} a_n|
\]
be two $n$-simplices containing $t$ as a face. Then the $(2n-1)$-coboundary satisfies
\[
\delta^{2n-1}(V^{t,s'}) = E^{s_1,s'} + E^{s_2,s'}.
\]
Hence $E^{s_1,s'}$ and $E^{s_2,s'}$ are cohomologous in $H^{2n}(\widetilde{K}^*;\mathbb{Z}_2)$. By the symmetry of $K$, this relation implies that for any pair of disjoint $n$-simplices $\sigma = a*\tau$ and $\sigma' = b*\tau'$ (where $\tau,\tau'\in\varDelta^{n-1}(\sigma_{2n}^{n-1})$ are disjoint), all such $E^{\sigma,\sigma'}$ are cohomologous to each other.


Next, fix an $(n-1)$-simplex $t' = |a_0 a_1 \cdots a_{n-1}|$ and an $n$-simplex $s' = |b a_{n+1} \cdots a_{2n}|$ in $K$. Let 
\[
s_1' = |a a_0 \cdots a_{n-1}|, \quad s_2' = |c a_0 \cdots a_{n-1}|
\]
be two $n$-simplices containing $t'$ as a face. Then, similarly,
\[
\delta^{2n-1}(V^{t',s'}) = E^{s_1',s'} + E^{s_2',s'}.
\]
By the symmetry of $K$, all $E^{\sigma,\sigma'}$ with $\sigma = a*\tau$ and $\sigma' = c*\tau'$ (where $\tau,\tau'\in\varDelta^{n-1}(\sigma_{2n}^{n-1})$ are disjoint) are also cohomologous to each other. The same type of relation holds for pairs of the form $\sigma = b*\tau$ and $\sigma' = c*\tau'$.


Combining these relations, we conclude that all dual cochains $E^{\sigma,\sigma'}$ corresponding to pairs of disjoint $n$-simplices in $K$ are cohomologous to each other. Therefore,
\[
H^{2n}(\widetilde{K}^*;\mathbb{Z}_2) \cong \mathbb{Z}_2,
\]
and this group is generated by the cohomology class $[c]$ of the dual cochain of any single $2n$-cell $\sigma \times \sigma'$, where $\sigma = a*\tau$ and $\sigma' = b*\tau'$ for disjoint $(n-1)$-simplices $\tau,\tau'\in\varDelta^{n-1}(\sigma_{2n}^{n-1})$. 
\end{proof}



\begin{Theorem}\label{vko}
For every generic immersion $\varphi$ of $K = \sigma_{2n}^{n-1}*\{a,b,c\}$ into $\mathbb{R}^{2n}$, the following holds:
\begin{eqnarray*}
\sum_{\substack{s,s'\in\varDelta^{n}(K)\\ s\cap s'=\varnothing}} l(\varphi(s),\varphi(s')) \equiv 1 \pmod{2}.
\end{eqnarray*}
\end{Theorem}

\begin{proof}
Let $J$ be a simplicial $(n-1)$-complex and let $a$, $b$, $c$ be three vertices not belonging to $J$. It is known that if the mod $2$ van Kampen obstruction $o_J$ of $J$ is non-zero, then the obstruction $o_{J*\{a,b,c\}}$ of the join $J*\{a,b,c\}$ is also non-zero \cite{Pa20}. Since the mod $2$ van Kampen obstruction $o_{\sigma_{2n}^{n-1}}$ is known to be non-zero, it follows that the obstruction $o_K$ of $K=\sigma_{2n}^{n-1}*\{a,b,c\}$ is also non-zero. 

On the other hand, by Lemma \ref{cohomo}, we have $H^{2n}(\widetilde{K}^*;\mathbb{Z}_2)\cong\mathbb{Z}_2$ and $o_K$ coincides with its unique non-trivial generator. Since all $2n$-cells of $\widetilde{K}^*$ are cohomologous, for every generic immersion $\varphi \colon K \to \mathbb{R}^{2n}$, we have
\begin{eqnarray*}
o_K = \bigg( \sum_{\substack{s,s'\in\varDelta^{n}(K)\\ s\cap s'=\emptyset}} l(\varphi(s),\varphi(s')) \bigg) \cdot [E^{\sigma,\sigma'}] \in H^{2n}(\widetilde{K}^*;\mathbb{Z}_2),
\end{eqnarray*}
where $\sigma = a*\tau$ and $\sigma' = b*\tau'$ for some disjoint $(n-1)$-simplices $\tau,\tau'$. Since $o_K$ is non-zero, the coefficient must be $1$ in $\mathbb{Z}_2$. This implies that the total number of double points is odd.
\end{proof}


\begin{thebibliography}{99}
\bibitem[CG83]{CG83} J. H. Conway and C. McA. Gordon, Knots and links in spatial graphs, \textit{J. Graph Theory} {\bf 7} (1983), 445--453.
\bibitem[Flores32]{Flores32} A. Flores, \"{U}ber $n$-dimensionale Komplexe die im $R_{2n+1}$ absolut selbstverschlungen sind, \textit{Ergeb. Math. Kolloq.} {\bf 6} (1932/1934), 4--7.
\bibitem[HRJV26]{HRJV26} N. Huber, I. R. Rao, H. Schwartz and T. T. Vijay, Intrinsic linking of $2$-complexes in $\mathbb{R}^{4}$, arXiv:math.GT/2605.06851.
\bibitem[N26b]{N26b} R. Nikkuni, Intrinsic linking of a simplicial $n$-complex embedded in $\mathbb{R}^{2n}$, arXiv:math.GT/2602.19511.
\bibitem[Pa20]{Pa20} S. Parsa, On the Smith classes, the van Kampen obstruction and embeddability of $[3] * K$, arXiv:math.AT/2001.06478.
\bibitem[VK33]{VK33} E. R. van Kampen, Komplexe in euklidischen R\"{a}umen, \textit{Abh. Math. Seminar Univ. Hamburg} {\bf 9} (1933), 72--78.
\end{thebibliography}
\end{document}
